\documentclass[10pt, aspectratio=169]{beamer}
% Preámbulo modular para presentaciones Beamer (Modelación Estadística)
% Diseño y paleta institucional del Tecnológico de Monterrey

\documentclass[aspectratio=169,xcolor={dvipsnames,table}]{beamer}

\usepackage[utf8]{inputenc}
\usepackage[spanish,mexico]{babel}
\usepackage{amsmath,amssymb,amsfonts,mathtools}
\usepackage{graphicx}
\usepackage{listings}
\usepackage{booktabs}
\usepackage{tikz}
\usetikzlibrary{matrix,arrows,decorations.pathmorphing,shapes.geometric,positioning}

% Paleta Institucional Tecnológico de Monterrey
\definecolor{TecAzul}{HTML}{0039A6}       % Azul TEC: color primario institucional
\definecolor{TecAzulOscuro}{HTML}{1A2E51} % Azul marino oscuro
\definecolor{TecRojo}{HTML}{EC2661}       % Rojo de acento (paleta miTec)
\definecolor{TecGris}{HTML}{646464}       % Gris neutro para comentarios y subtítulos
\definecolor{TecFondoBloque}{HTML}{F4F6F9}% Fondo suave para bloques

% Configuración de Tema Beamer
\usetheme{Boadilla}
\usecolortheme[named=TecAzulOscuro]{structure}
\setbeamercolor{alerted text}{fg=TecRojo}
\setbeamercolor{palette primary}{bg=TecAzulOscuro,fg=white}
\setbeamercolor{palette secondary}{bg=TecAzul,fg=white}
\setbeamercolor{palette tertiary}{bg=TecAzulOscuro!80!black,fg=white}
\setbeamercolor{title}{fg=TecAzulOscuro,bg=TecFondoBloque}
\setbeamercolor{frametitle}{fg=TecAzulOscuro,bg=TecFondoBloque}

% Bloques personalizados elegantes
\setbeamercolor{block title}{bg=TecAzulOscuro,fg=white}
\setbeamercolor{block body}{bg=TecFondoBloque,fg=black}
\setbeamercolor{block title alerted}{bg=TecRojo,fg=white}
\setbeamercolor{block body alerted}{bg=TecRojo!8,fg=black}
\setbeamercolor{block title example}{bg=TecAzul,fg=white}
\setbeamercolor{block body example}{bg=TecAzul!8,fg=black}

% Redondear bordes y añadir sombra sutil
\setbeamertemplate{blocks}[rounded][shadow=true]
\setbeamertemplate{navigation symbols}{}

% Pie de página limpio con número de diapositiva
\setbeamertemplate{footline}{
  \leavevmode%
  \hbox{%
  \begin{beamercolorbox}[wd=.7\paperwidth,ht=2.25ex,dp=1ex,left]{author in head/foot}%
    \hspace*{2ex}\insertshortauthor~~(\insertshortinstitute) \hfill \insertshorttitle
  \end{beamercolorbox}%
  \begin{beamercolorbox}[wd=.3\paperwidth,ht=2.25ex,dp=1ex,right]{date in head/foot}%
    \insertshortdate{}\hspace*{2em}
    \insertframenumber{} / \inserttotalframenumber\hspace*{2ex}
  \end{beamercolorbox}}%
  \vskip0pt%
}

% Configuración de Listings de Python para visualización pedagógica de código
\lstset{
    language=Python,
    basicstyle=\ttfamily\footnotesize,
    keywordstyle=\color{TecAzulOscuro}\bfseries,
    stringstyle=\color{TecRojo},
    commentstyle=\color{TecGris}\itshape,
    showstringspaces=false,
    breaklines=true,
    frame=lines,
    rulecolor=\color{TecAzulOscuro!30},
    backgroundcolor=\color{TecFondoBloque},
    aboveskip=0.5em,
    belowskip=0.5em,
    literate={á}{{\'a}}1 {é}{{\'e}}1 {í}{{\'i}}1 {ó}{{\'o}}1 {ú}{{\'u}}1
             {Á}{{\'A}}1 {É}{{\'E}}1 {Í}{{\'I}}1 {Ó}{{\'O}}1 {Ú}{{\'U}}1
             {ñ}{{\~n}}1 {Ñ}{{\~N}}1 {¿}{{?`}}1 {¡}{{!`}}1
}

% Metadatos institucionales por defecto
\title[Modelación Estadística]{Modelación Estadística para Ciencia de Datos e Ingeniería}
\author[J. Castillo Colmenares]{Juliho David Castillo Colmenares}
\institute[Tec de Monterrey]{Tecnológico de Monterrey \\ Escuela de Ingeniería y Ciencias}
\date{\today}

% Comandos y macros estadísticas compartidas para las presentaciones Beamer (Metropolis)

% Conjuntos numéricos básicos
\newcommand{\R}{\mathbb{R}}
\newcommand{\N}{\mathbb{N}}
\newcommand{\Q}{\mathbb{Q}}
\newcommand{\Z}{\mathbb{Z}}
\newcommand{\set}[1]{\left\{ #1 \right\}}
\newcommand{\abs}[1]{\left|#1\right|}
\newcommand{\norm}[1]{\left\| #1 \right\|}

% Comandos estadísticos y probabilísticos del curso
\newcommand{\Var}{\operatorname{Var}}
\newcommand{\cov}{\operatorname{Cov}}
\newcommand{\corr}{\rho}
\newcommand{\comb}[2]{\begin{pmatrix} #1 \\ #2 \end{pmatrix}}
\newcommand{\s}{\sigma}
\newcommand{\particion}{\mathfrak{P}}
\newcommand{\card}[1]{\#\left( #1 \right)}
\newcommand{\seq}[3]{\set{#1}_{#2}^{#3}}
\newcommand{\E}{\mathbb{E}}
\newcommand{\Prob}{\mathbb{P}}

% Entornos pedagógicos y cajas en español académico natural
\newenvironment{discusion}{%
  \begin{alertblock}{Pregunta de Discusión}%
}{%
  \end{alertblock}%
}

\newenvironment{reflexion}{%
  \begin{alertblock}{Reflexión para el Aula}%
}{%
  \end{alertblock}%
}

\newenvironment{paradoja}{%
  \begin{alertblock}{Paradoja de Exploración}%
}{%
  \end{alertblock}%
}

\newenvironment{motivacion}{%
  \begin{exampleblock}{Motivación: Ciencia de Datos en la Práctica}%
}{%
  \end{exampleblock}%
}

\newenvironment{retoclase}{%
  \begin{block}{Reto Rápido en Equipo}%
}{%
  \end{block}%
}


\title[Pruebas $Z$/$t$ para Medias]{Sección 07.02: Pruebas $Z$ y $t$ para Medias de Una y Dos Muestras}
\subtitle{Independientes, Homoscedásticas, de Welch y Pareadas}
\author[J. Castillo Colmenares]{Juliho Castillo Colmenares}
\institute{Tecnológico de Monterrey}
\date{\vspace{-1.2cm}}

\begin{document}

% Slide 1: Portada
\begin{frame}[plain]
  \titlepage
\end{frame}

% Slide 2: Hoja de Ruta
\begin{frame}{Hoja de Ruta --- Capítulo 07: Docimasia (Pruebas de Hipótesis)}
  \footnotesize
  ¿Dónde nos encontramos en el desarrollo modular del capítulo? \pause
  \begin{itemize}\itemsep=0.08em
    \item \textbf{07.01} Fundamentos: $H_0$ vs. $H_1$, Errores Tipo I y II, y Potencia \pause
    \item \textbf{07.02} Pruebas $Z$ y $t$ para Medias de Una y Dos Muestras \emph{(Hoy)} \pause
    \item \textbf{07.03} Pruebas de Bondad de Ajuste $\chi^2$ \pause
    \item \textbf{07.04} Tablas de Contingencia y Pruebas de Independencia
  \end{itemize}
  \vspace{-0.05cm}
  \begin{block}{Objetivo de la Sesión}
    Extender el marco lógico de la Sección 07.01 (planteamiento de hipótesis, $\alpha$, $\beta$, potencia) a las pruebas concretas para una y dos medias: la prueba $t$ de una muestra, el caso de dos muestras independientes homoscedásticas (varianza agrupada), el caso heteroscedástico (Welch), y el caso de muestras pareadas.
  \end{block}
\end{frame}

% Slide 3: Motivación
\begin{frame}{De una Media a la Comparación de Grupos}
  \footnotesize
  \begin{motivacion}
    En la Sección 07.01 aplicamos el marco de docimasia a una sola media con $\sigma$ conocida. En la práctica, casi nunca conocemos $\sigma$, y con frecuencia el objetivo real es \emph{comparar} dos grupos: ¿es el algoritmo A más rápido que B? ¿mejora el tratamiento el desempeño de los mismos sujetos?
  \end{motivacion}

  \vspace{0.15cm}
  \pause
  \begin{itemize}\itemsep=0.1cm
    \item \textbf{$\sigma$ desconocida (una muestra):} reemplazamos $\sigma$ por $S$ y usamos la distribución $t$ con $n-1$ g.l. \pause
    \item \textbf{Dos muestras independientes:} el estadístico depende de si las varianzas poblacionales son iguales o distintas. \pause
    \item \textbf{Muestras pareadas:} se reducen algebraicamente a una prueba $t$ de una sola muestra sobre las diferencias.
  \end{itemize}
\end{frame}

% Slide 4: Prueba t de una muestra
\begin{frame}{Prueba $t$ de Una Muestra ($\sigma$ Desconocida)}
  \footnotesize
  \begin{block}{Estadístico de Prueba}
    $$t = \frac{\bar X - \mu_0}{S/\sqrt n} \ \sim\ t_{n-1} \quad \text{bajo } H_0: \mu=\mu_0$$
  \end{block}
  \pause

  \vspace{0.1cm}
  \begin{alertblock}{Comparación con la Sección 07.01}
    La estructura es idéntica a la prueba $Z$, salvo que: (1) $\sigma$ se reemplaza por su estimador muestral $S$, y (2) la distribución de referencia pasa de $N(0,1)$ a $t_{n-1}$, más dispersa para compensar la incertidumbre adicional de estimar $\sigma$.
  \end{alertblock}
\end{frame}

% Slide 5: Dos medias independientes (pooled y Welch)
\begin{frame}{Dos Medias Independientes: Varianza Agrupada vs. Welch}
  \footnotesize
  \begin{block}{Caso Homoscedástico ($\sigma_1^2=\sigma_2^2$): Varianza Agrupada}
    $$t = \frac{(\bar X_1-\bar X_2)-\Delta_0}{S_p\sqrt{1/n_1+1/n_2}}, \qquad S_p^2 = \frac{(n_1-1)S_1^2+(n_2-1)S_2^2}{n_1+n_2-2}, \qquad \nu=n_1+n_2-2$$
  \end{block}
  \pause

  \vspace{0.1cm}
  \begin{alertblock}{Caso Heteroscedástico ($\sigma_1^2\neq\sigma_2^2$): Welch-Satterthwaite}
    $$t = \frac{(\bar X_1-\bar X_2)-\Delta_0}{\sqrt{S_1^2/n_1+S_2^2/n_2}}, \qquad \nu_{\text{Welch}} = \frac{(S_1^2/n_1+S_2^2/n_2)^2}{\frac{(S_1^2/n_1)^2}{n_1-1}+\frac{(S_2^2/n_2)^2}{n_2-1}}$$
  \end{alertblock}
\end{frame}

% Slide 6: Muestras pareadas
\begin{frame}{Muestras Pareadas: Reducción a una Sola Muestra}
  \footnotesize
  \begin{block}{Variable de Diferencia}
    Para mediciones repetidas sobre los mismos sujetos ($X_{1i}, X_{2i}$), definimos $D_i=X_{1i}-X_{2i}$ y reducimos el contraste sobre $\mu_1-\mu_2$ a una prueba $t$ de una sola muestra sobre $\mu_D$:
    $$t = \frac{\bar D - \Delta_0}{S_D/\sqrt n} \ \sim\ t_{n-1}$$
  \end{block}
  \pause

  \vspace{0.1cm}
  \begin{alertblock}{¿Por Qué Pareamos?}
    El pareamiento elimina la variabilidad \emph{entre sujetos} (cada quien es su propio control), aumentando drásticamente la potencia de la prueba al reducir el error estándar frente al diseño de muestras independientes.
  \end{alertblock}
\end{frame}

% Slide 7: Elegir la prueba correcta
\begin{frame}{¿Qué Prueba Elegir? Un Árbol de Decisión}
  \footnotesize
  \begin{itemize}\itemsep=0.12cm
    \item \textbf{¿Una sola muestra?} Use la prueba $t$ (o $Z$ si $\sigma$ es conocida) de la Sección 07.01. \pause
    \item \textbf{¿Dos muestras, mismos sujetos medidos dos veces?} Use la prueba $t$ pareada sobre $D_i$. \pause
    \item \textbf{¿Dos muestras independientes, varianzas similares?} Use la prueba $t$ con varianza agrupada. \pause
    \item \textbf{¿Dos muestras independientes, varianzas claramente distintas?} Use la prueba $t$ de Welch.
  \end{itemize}
\end{frame}

% Slide 8: Lab Python (Bloque 1)
\begin{frame}[fragile]{Laboratorio en Python: Prueba $t$ de Una Muestra}
  \scriptsize
  Prueba $t$ con $\sigma$ desconocida, verificada contra \texttt{scipy.stats.ttest\_1samp} (\texttt{07.02\_z\_t\_tests\_means.py}):
  {\lstset{basicstyle=\fontsize{4pt}{4.8pt}\selectfont\ttfamily,aboveskip=0.1em,belowskip=0.1em}
  \lstinputlisting[firstline=17, lastline=37]{../../code/07_pruebas_hipotesis/07.02_z_t_tests_means.py}}
\end{frame}

% Slide 9: Lab Python (Bloque 2)
\begin{frame}[fragile]{Laboratorio en Python: Varianza Agrupada vs. Welch}
  \scriptsize
  Comparación directa de ambos estadísticos bajo homoscedasticidad y heteroscedasticidad:
  {\lstset{basicstyle=\fontsize{3.6pt}{4.3pt}\selectfont\ttfamily,aboveskip=0.05em,belowskip=0.05em}
  \lstinputlisting[firstline=40, lastline=77]{../../code/07_pruebas_hipotesis/07.02_z_t_tests_means.py}}
\end{frame}

% Slide 10: Lab Python (Bloque 3)
\begin{frame}[fragile]{Laboratorio en Python: Prueba $t$ Pareada}
  \scriptsize
  Reducción de un diseño pareado a una prueba $t$ de una sola muestra sobre las diferencias:
  {\lstset{basicstyle=\fontsize{4pt}{4.8pt}\selectfont\ttfamily,aboveskip=0.1em,belowskip=0.1em}
  \lstinputlisting[firstline=80, lastline=101]{../../code/07_pruebas_hipotesis/07.02_z_t_tests_means.py}}
\end{frame}

% Slide 11: Salida Terminal
\begin{frame}[fragile]{Laboratorio en Python: Salida en Terminal}
  \scriptsize
  Salida estándar generada al ejecutar el script del laboratorio en Python:
  \vspace{0.02cm}
  \begin{block}{Salida Estándar en Consola}
    \fontsize{4pt}{5pt}\selectfont
    \begin{verbatim}
=== Block 1: One-Sample t-Test ===
t statistic = 1.7050, df=48, two-tailed p-value = 0.0947

=== Block 2: Pooled vs. Welch t-Test ===
Homoscedastic: Pooled t=-1.9815, df=32, p=0.0562
Heteroscedastic: Welch t=4.2026, df=16.60, p=0.0006

=== Block 3: Paired t-Test ===
t statistic = 6.3743, df=11, one-tailed p-value = 0.000026
    \end{verbatim}
  \end{block}
\end{frame}

% Slide 12A: Ejercicio Nivel 1 (Enunciado)
\begin{frame}{Ejercicio en Clase --- Nivel 1: Fundamental (1/2)}
  \footnotesize
  \begin{block}{Problema (Enunciado, Docimasia Avanzada)}
    Una empresa de telemedicina compara el tiempo medio de respuesta entre dos plataformas: Plataforma 1 ($n_1=16$, $\bar X_1=14.2$, $S_1=2.4$) y Plataforma 2 ($n_2=18$, $\bar X_2=16.8$, $S_2=2.6$), asumiendo varianzas iguales. Pruebe al nivel $\alpha=0.05$ si existe diferencia real.
  \end{block}

  \vspace{0.15cm}
  \pause
  \begin{alertblock}{Planteamiento}
    Calcule primero la varianza agrupada $S_p^2$ con $\nu=n_1+n_2-2=32$ grados de libertad.
  \end{alertblock}
\end{frame}

% Slide 12B: Ejercicio Nivel 1 (Resolución)
\begin{frame}{Ejercicio en Clase --- Nivel 1: Fundamental (2/2)}
  \scriptsize
  Calculamos la varianza agrupada y el estadístico $t$: \pause

  \vspace{0.05cm}
  \begin{block}{Cálculo}
    $$S_p^2 = \frac{15(2.4)^2+17(2.6)^2}{32} \approx 6.291, \qquad t = \frac{14.2-16.8}{2.508\sqrt{1/16+1/18}} \approx -3.017$$
  \end{block}

  \vspace{0.05cm}
  \pause
  \begin{alertblock}{Interpretación}
    Con $t_{0.025,32}\approx2.037$, como $|-3.017|>2.037$ se \textbf{rechaza $H_0$}: la Plataforma 1 es significativamente más rápida al nivel del 5\%.
  \end{alertblock}
\end{frame}

% Slide 13A: Ejercicio Nivel 2 (Enunciado)
\begin{frame}{Ejercicio en Clase --- Nivel 2: Operativo (1/2)}
  \footnotesize
  \begin{block}{Problema (Enunciado, Docimasia Avanzada)}
    Un científico de datos compara tiempos de ejecución de dos algoritmos con varianzas explícitamente heterogéneas: Algoritmo 1 ($n_1=15$, $\bar X_1=85.2$, $S_1=12.4$), Algoritmo 2 ($n_2=18$, $\bar X_2=74.5$, $S_2=5.2$). Calcule $\nu$ de Welch y pruebe al $\alpha=0.05$.
  \end{block}

  \vspace{0.15cm}
  \pause
  \begin{alertblock}{Estrategia}
    \scriptsize
    Use la fórmula de Welch-Satterthwaite para los grados de libertad efectivos (generalmente no enteros).
  \end{alertblock}
\end{frame}

% Slide 13B: Ejercicio Nivel 2 (Resolución)
\begin{frame}{Ejercicio en Clase --- Nivel 2: Operativo (2/2)}
  \scriptsize
  Calculamos los grados de libertad de Welch y el estadístico $t$: \pause

  \vspace{0.05cm}
  \begin{block}{Cálculo}
    \scriptsize
    $$\nu_{\text{Welch}} \approx 18.14, \qquad t = \frac{85.2-74.5}{\sqrt{12.4^2/15+5.2^2/18}} \approx 3.14$$
  \end{block}

  \vspace{0.05cm}
  \pause
  \begin{alertblock}{Interpretación}
    \scriptsize
    Con $t_{0.025,18}\approx2.101$, como $3.14>2.101$ se \textbf{rechaza $H_0$}: los tiempos medios de ejecución difieren significativamente; usar la varianza agrupada aquí habría sido metodológicamente incorrecto.
  \end{alertblock}
\end{frame}

% Slide 14A: Ejercicio Nivel 3 (Enunciado)
\begin{frame}{Ejercicio en Clase --- Nivel 3: Analítico (1/2)}
  \footnotesize
  \begin{block}{Problema --- Potencia y Tamaño de Muestra (Enunciado, Docimasia Avanzada)}
    Para $H_0:\mu\le100$ vs. $H_a:\mu>100$ ms con $\sigma=15$ conocida, $n=25$, $\alpha=0.05$: calcule $\beta$ si $\mu_a=108$, y el $n$ necesario para potencia del $90\%$.
  \end{block}

  \vspace{0.1cm}
  \pause
  \begin{alertblock}{Estrategia}
    \scriptsize
    Este problema conecta directamente con la fórmula de tamaño de muestra de la Sección 07.01: $n=((Z_\alpha+Z_\beta)\sigma/(\mu_a-\mu_0))^2$.
  \end{alertblock}
\end{frame}

% Slide 14B: Ejercicio Nivel 3 (Resolución)
\begin{frame}{Ejercicio en Clase --- Nivel 3: Analítico (2/2)}
  \scriptsize
  Aplicamos las fórmulas de la Sección 07.01 al contexto de latencia de red: \pause

  \vspace{0.05cm}
  \begin{block}{Cálculo}
    \scriptsize
    $$\beta \approx 0.153 \implies \text{Potencia}\approx0.847, \qquad n = \left(\frac{(1.645+1.282)(15)}{8}\right)^2 \approx 30.12 \to n=31$$
  \end{block}

  \vspace{0.05cm}
  \pause
  \begin{alertblock}{Interpretación}
    \scriptsize
    Este ejercicio confirma que la maquinaria teórica de la Sección 07.01 (diseñada para una sola media con $\sigma$ conocida) se aplica sin cambios a cualquier contexto aplicado, aquí latencia de red en una fintech.
  \end{alertblock}
\end{frame}

% Slide 15A: Ejercicio Nivel 4 (Enunciado)
\begin{frame}{Ejercicio en Clase --- Nivel 4: Desafiante (1/2)}
  \footnotesize
  \begin{block}{Problema --- Razón de Verosimilitudes de Wilks (Enunciado, Docimasia Avanzada)}
    Para $X_i\sim N(\mu_1,\sigma_1^2)$, $Y_i\sim N(\mu_2,\sigma_2^2)$ independientes, deduzca el estadístico LRT $-2\ln\Lambda$ para $H_0:\sigma_1^2=\sigma_2^2$ y justifique su convergencia a $\chi^2_1$ vía el Teorema de Wilks.
  \end{block}

  \vspace{0.1cm}
  \pause
  \begin{alertblock}{Estrategia}
    \scriptsize
    Derive primero los EMV de $\sigma_1^2,\sigma_2^2$ sin restricción, luego el EMV combinado bajo la restricción $\sigma_1^2=\sigma_2^2$.
  \end{alertblock}
\end{frame}

% Slide 15B: Ejercicio Nivel 4 (Resolución)
\begin{frame}{Ejercicio en Clase --- Nivel 4: Desafiante (2/2)}
  \scriptsize
  Construimos el estadístico de razón de verosimilitudes: \pause

  \vspace{0.05cm}
  \begin{block}{Cálculo}
    \scriptsize
    $$\hat\sigma_0^2 = \frac{n_1\hat\sigma_1^2+n_2\hat\sigma_2^2}{n_1+n_2}, \qquad -2\ln\Lambda = n_1\ln\left(\frac{\hat\sigma_0^2}{\hat\sigma_1^2}\right) + n_2\ln\left(\frac{\hat\sigma_0^2}{\hat\sigma_2^2}\right)$$
  \end{block}

  \vspace{0.05cm}
  \pause
  \begin{alertblock}{Interpretación}
    \scriptsize
    Por el Teorema de Wilks, $-2\ln\Lambda \to \chi^2_{\nu}$ donde $\nu$ es la diferencia de dimensión paramétrica entre $\Theta$ (2 parámetros de varianza libres) y $\Theta_0$ (1 parámetro combinado), es decir $\nu=1$. Este es un enfoque alternativo y asintóticamente equivalente a la prueba $F$ clásica.
  \end{alertblock}
\end{frame}

% Slide 16: Cierre
\begin{frame}{Cierre de Sección: Pruebas $Z$/$t$ para Medias}
  \begin{reflexion}
    Comparar medias exige primero diagnosticar la estructura del problema --- ¿una muestra o dos? ¿independientes o pareadas? ¿varianzas iguales o distintas? --- antes de aplicar mecánicamente una fórmula. Elegir el estadístico incorrecto (por ejemplo, varianza agrupada cuando las varianzas son claramente distintas) puede invalidar la inferencia.
  \end{reflexion}

  \vspace{0.2cm}
  \pause
  \begin{block}{Lecciones Clave}
    \begin{itemize}\itemsep=0.08cm
      \item \textbf{Una muestra, $\sigma$ desconocida:} $t=(\bar X-\mu_0)/(S/\sqrt n)$, $\nu=n-1$.
      \item \textbf{Dos muestras independientes:} varianza agrupada si $\sigma_1^2=\sigma_2^2$; Welch si no.
      \item \textbf{Muestras pareadas:} reducir a una prueba $t$ de una muestra sobre las diferencias $D_i$.
    \end{itemize}
  \end{block}
\end{frame}

% Slide 17: Perspectiva Modular
\begin{frame}{Perspectiva Modular --- El Próximo Reto}
  \scriptsize
  \begin{retoclase}
    Hasta ahora hemos probado hipótesis sobre medias de variables \emph{continuas}. La Sección 07.03 cambia de terreno: probaremos si los datos \emph{categóricos} completos se ajustan a una distribución teórica propuesta, usando la prueba de bondad de ajuste $\chi^2$ --- la herramienta fundamental para validar modelos discretos completos, no solo un parámetro aislado.
  \end{retoclase}

  \vspace{0.08cm}
  \pause
  \begin{alertblock}{Preparación Recomendada}
    \begin{itemize}\itemsep=0.06cm
      \item Repasa la distribución $\chi^2$ y sus grados de libertad (Sección 05.03).
      \item Reflexiona sobre por qué estimar parámetros de la propia muestra reduce los grados de libertad disponibles.
    \end{itemize}
  \end{alertblock}
\end{frame}

\end{document}
